\chapter{From Particle Mechanics to Classical Fields}

\section{What a Field Is}

The preceding lecture developed relativistic particle mechanics from an invariant
action.  We now change the dynamical object.  The subject is classical field
theory, not quantum field theory.  Quantum mechanics will occasionally provide
a useful comparison, but the equations in this chapter describe ordinary
classical fields.  Electromagnetism is the familiar model in the background:
electromagnetic waves propagate through space and interact with charged
particles.  We begin with a still simpler field.

A spacetime may have many spatial coordinates, but it has one time coordinate.
Write
\begin{equation}
x^\mu=(t,x^1,\ldots,x^d),
\end{equation}
where the spatial index \(i=1,\ldots,d\).  The number \(d\) may be three, as in
our world, or one, two, nine, twenty-five, or any other number useful in a
model.  The important asymmetry at this stage is that there is still only one
time direction.  The lecture does not try to explain why theories with several
times are so difficult; it simply adopts the one-time framework in which
ordinary causal physics is formulated.

A field is a measurable quantity whose value may vary from point to point and
from time to time.  Temperature is a simple example:
\begin{equation}
T=T(\mathbf{x},t).
\end{equation}
It assigns one number to every event, so it is a scalar field.  Wind velocity
also varies over space and time, but it has magnitude and direction and
therefore several components; it is a vector field.  In relativistic physics
the natural vectors have four components rather than three.

For most of this lecture the field has one component and is denoted by the
Greek letter corresponding to \(f\):
\begin{equation}
\phi=\phi(\mathbf{x},t).
\label{eq:sr4-field-definition}
\end{equation}
Calling \(\phi\) a scalar will eventually mean more than saying that it has one
component.  It will mean that all inertial observers assign the same value to
\(\phi\) at the same physical event.  That transformation law enters only
after the variational mechanics has been built.

\section{The Particle Prototype}

\subsection{A coordinate called \texorpdfstring{\(\phi\)}{phi}}

The quickest route to field theory begins by stepping back to an ordinary
nonrelativistic particle.  Let the particle move on a line, but call its
coordinate \(\phi\) instead of \(x\).  There is no mathematical obligation to
name a coordinate \(x\); the unusual name is chosen to make the coming analogy
visible.  Draw time horizontally.  The entire history of the particle is then
a curve \(\phi(t)\), which may pass through positive and negative values.

For fixed initial and final times \(a\) and \(b\), the action is
\begin{align}
A[\phi]
&=\int_a^b dt\,L(\phi,\dot\phi),\\
L(\phi,\dot\phi)
&=\frac{m}{2}\dot\phi^{\,2}-V(\phi).
\label{eq:sr4-particle-lagrangian}
\end{align}
The Euler--Lagrange equation is
\begin{equation}
\frac{d}{dt}\left(\frac{\partial L}{\partial\dot\phi}\right)
-\frac{\partial L}{\partial\phi}=0.
\label{eq:sr4-particle-el}
\end{equation}
For the Lagrangian in \eqref{eq:sr4-particle-lagrangian},
\begin{equation}
\frac{\partial L}{\partial\dot\phi}=m\dot\phi,
\qquad
\frac{\partial L}{\partial\phi}=-\frac{dV}{d\phi},
\end{equation}
so
\begin{equation}
\boxed{m\ddot\phi=-\frac{dV}{d\phi}}.
\label{eq:sr4-newton}
\end{equation}
This is simply Newton's equation, written with \(\phi\) as the coordinate.

\begin{figure}[htbp]
\centering
\includegraphics[width=0.90\textwidth]{lecture_04_figure_01.png}
\caption{The ordinary Euler--Lagrange equation and its Newtonian result for the coordinate \(\phi(t)\).}
\lectureframe{00:10:18}
\end{figure}

The physical path makes the action stationary while the endpoints are held
fixed.  The lecture often says \emph{minimum}, but the precise statement is
stationarity: the path may give a minimum, a maximum, or another stationary
value.  One useful mental model is to divide the time interval into many small
pieces.  The action becomes a function of the many sampled values of \(\phi\).
Moving those values, while leaving the endpoints fixed, and differentiating
with respect to each of them is the finite-dimensional ancestor of the
Euler--Lagrange equation.

\subsection{Mechanics as a theory in \texorpdfstring{\(0+1\)}{0+1} dimensions}

Now imagine a world with time but no spatial dimensions.  Its space is a
single point.  A field in that world cannot depend on position, because there
is no position on which to depend; it can only be a function \(\phi(t)\).
Ordinary one-coordinate particle mechanics can therefore be read as a field
theory in \(0+1\) dimensions.  The mathematics has not changed.  Only the
interpretation has changed from the position of a particle to the value of a
field in a world with no space.

\begin{figure}[htbp]
\centering
\includegraphics[width=0.90\textwidth]{lecture_04_figure_02.png}
\caption{The sampled history \(\phi(t)\): a particle trajectory and, with a new interpretation, a field history in \(0+1\) dimensions.}
\lectureframe{00:14:30}
\end{figure}

\begin{classroomqa}
\audiencequestion{Does the potential disappear when there are no spatial dimensions?}
\lecturerresponse{No.  Spatial variation disappears, but the field still has a value.  The energy may depend on that value through \(V(\phi)\).  In the particle interpretation, moving the curve up or down changes the potential energy; in the field interpretation, different values of \(\phi\) carry different potential energy. \lecturetimestamp{00:15:00}}
\end{classroomqa}

\begin{classroomqa}
\audiencequestion{Instead of imagining no space at all, may one imagine an ordinary space in which the field depends only on time?}
\lecturerresponse{Yes.  For this purpose the local mathematics is the same: all spatial derivatives vanish and only \(\phi(t)\) remains.  Calling it a \(0+1\)-dimensional theory isolates that reduced structure, but a spatially homogeneous field in a larger world supplies the same intuition. \lecturetimestamp{00:16:25}}
\end{classroomqa}

The point of the detour is now clear.  We know the action principle for a field
that depends on one coordinate, time.  We can use its pattern to guess the
action principle for a field that depends on all the coordinates of spacetime.

\section{A Field History over Spacetime}

\subsection{From a curve to a surface}

In \(1+1\) dimensions, spacetime itself is a plane with coordinates \(t\) and
\(x\).  The value \(\phi(x,t)\) supplies a third direction.  The history of the
field is therefore not a curve but a surface above the spacetime plane.  The
surface may rise above or fall below zero, because a scalar field need not be
positive.

\begin{figure}[htbp]
\centering
\includegraphics[width=0.90\textwidth]{lecture_04_figure_03.png}
\caption{A scalar-field history drawn as a surface over a \(1+1\)-dimensional spacetime plane.}
\lectureframe{00:20:46}
\end{figure}

Our physical spacetime is called \(3+1\)-dimensional: three dimensions of space
and one of time.  A tabletop model is \(2+1\)-dimensional, and beads constrained
to a wire suggest a \(1+1\)-dimensional world.  The extra axis needed to plot
\(\phi\) makes even the \(3+1\) case impossible to draw faithfully, but the
definition is unambiguous: a complete history assigns one value of \(\phi\) to
every spacetime event.

\subsection{From fixed endpoints to fixed boundary data}

The particle variational problem fixes two endpoints and varies the curve
between them.  Its field-theory analog selects a spacetime region, fixes
\(\phi\) on the boundary of that region, and varies \(\phi\) in the interior.
The region may be a rectangle in a \(1+1\)-dimensional drawing or any suitable
region in a higher-dimensional spacetime.

\begin{figure}[htbp]
\centering
\includegraphics[width=0.90\textwidth]{lecture_04_figure_04.png}
\caption{A spacetime region whose entire boundary is held fixed while the field in the interior is varied.}
\lectureframe{00:24:52}
\end{figure}

\begin{classroomqa}
\audiencequestion{Does the Lagrangian formulation still work when a potential depends explicitly on time?}
\lecturerresponse{Yes.  The Euler--Lagrange and Hamiltonian formulations still work.  What is lost, in general, is energy conservation: explicit time dependence breaks time-translation symmetry.  The corresponding statement in field theory is that explicit spacetime dependence is allowed, although it affects the associated conservation laws. \lecturetimestamp{00:23:49}}
\end{classroomqa}

\begin{classroomqa}
\audiencequestion{For the rectangular field region, is only the top and bottom held fixed, or the whole boundary?}
\lecturerresponse{For this variational construction, the whole boundary is fixed: all four sides in the \(1+1\)-dimensional rectangle.  The field is allowed to wiggle only in the interior.  This is the direct analog of holding both endpoints fixed in particle mechanics. \lecturetimestamp{00:24:24}}
\end{classroomqa}

Fixing the whole boundary is a statement about the variation used to derive the
equations.  When the resulting hyperbolic field equation is treated as a
physical evolution problem, one will more commonly specify appropriate initial
data and let the equation propagate it forward.

\section{The Action and the Field Euler--Lagrange Equation}

\subsection{The Lagrangian density}

Divide the spacetime region into small cells.  The action is assumed to be a
sum of local contributions, one from each cell.  In the continuum limit that
sum becomes an integral over spacetime:
\begin{equation}
A[\phi]
=
\int_{\Omega} d^4x\,
\mathcal{L}\!\left(
\phi,\partial_\mu\phi;x^\mu
\right),
\qquad
d^4x=dt\,dx\,dy\,dz.
\label{eq:sr4-field-action}
\end{equation}
The quantity \(\mathcal{L}\) is the Lagrangian density.  The semicolon in
\(\mathcal{L}(\phi,\partial_\mu\phi;x^\mu)\) separates dynamical arguments from
possible explicit coordinate dependence.  In a homogeneous closed system that
last dependence is normally absent, but the variational formalism does not
forbid it.

The compact notation
\begin{equation}
\partial_\mu\phi
\equiv
\frac{\partial\phi}{\partial x^\mu}
\end{equation}
collects the time and spatial derivatives.  In the lecture's component
notation one may instead write
\begin{equation}
\begin{aligned}
\phi_t&=\frac{\partial\phi}{\partial t},\\
\phi_x&=\frac{\partial\phi}{\partial x},\\
\phi_y&=\frac{\partial\phi}{\partial y},\\
\phi_z&=\frac{\partial\phi}{\partial z}.
\end{aligned}
\label{eq:sr4-component-derivatives}
\end{equation}
There may be several fields, just as a particle may have several coordinates.
For now there is only one field \(\phi\).

The symbol \(d^4x\) is only shorthand for the four-coordinate measure.  In a
spacetime of another dimension it would be replaced by \(d^Dx\).  In the
\(0+1\)-dimensional prototype, \(D=1\) and \(d^1x\) is simply \(dt\).

\begin{figure}[htbp]
\centering
\includegraphics[width=0.90\textwidth]{lecture_04_figure_05.png}
\caption{The action density written as a local function of \(\phi\) and its spacetime derivatives.}
\lectureframe{00:32:51}
\end{figure}

\begin{classroomqa}
\audiencequestion{May the Lagrangian density also depend on \(x,y,z\) themselves?}
\lecturerresponse{It may, just as a particle Lagrangian may depend explicitly on \(t\).  For a translation-invariant closed system it ordinarily does not, but explicit coordinate dependence does not invalidate the formalism. \lecturetimestamp{00:30:29}}
\end{classroomqa}

\begin{classroomqa}
\audiencequestion{Is the construction so far specifically for a scalar field, or is it general?}
\lecturerresponse{So far the variational pattern is general.  Writing only one component suggests a scalar field.  A vector field or a collection of fields would carry additional component labels, and one would vary every component. \lecturetimestamp{00:31:00}}
\end{classroomqa}

\subsection{The lattice picture}

On a fine spacetime lattice, derivatives become finite differences between
neighboring field values.  The action is then an ordinary function of many
variables.  A grid with ten points in one direction and eight in another
already gives eighty field variables.  Refining the grid increases their
number, but the minimization rule remains elementary: differentiate the action
with respect to the field value at each interior point and set the result to
zero.

\begin{figure}[htbp]
\centering
\includegraphics[width=0.90\textwidth]{lecture_04_figure_06.png}
\caption{The spacetime region divided into cells, making the action an ordinary many-variable function before the continuum limit.}
\lectureframe{00:34:05}
\end{figure}

Taking the continuum limit gives the field Euler--Lagrange equation,
\begin{equation}
\boxed{
\partial_\mu
\left[
\frac{\partial\mathcal{L}}
{\partial(\partial_\mu\phi)}
\right]
-
\frac{\partial\mathcal{L}}{\partial\phi}
=0
}.
\label{eq:sr4-field-el}
\end{equation}
Repeated \(\mu\) is summed over all spacetime directions.  Written without the
summation shorthand, the same equation is
\begin{align}
&
\frac{\partial}{\partial t}
\left(\frac{\partial\mathcal{L}}{\partial\phi_t}\right)
+
\frac{\partial}{\partial x}
\left(\frac{\partial\mathcal{L}}{\partial\phi_x}\right)
+
\frac{\partial}{\partial y}
\left(\frac{\partial\mathcal{L}}{\partial\phi_y}\right)
\nonumber\\
&\hspace{3em}
+
\frac{\partial}{\partial z}
\left(\frac{\partial\mathcal{L}}{\partial\phi_z}\right)
-
\frac{\partial\mathcal{L}}{\partial\phi}
=0.
\label{eq:sr4-field-el-expanded}
\end{align}
The first term is the familiar particle term.  Field theory adds an analogous
term for every spatial direction.  The last term,
\(-\partial\mathcal{L}/\partial\phi\), is unchanged.

\begin{figure}[htbp]
\centering
\includegraphics[width=0.90\textwidth]{lecture_04_figure_07.png}
\caption{The one-dimensional Euler--Lagrange equation and its spacetime generalization, written one below the other.}
\lectureframe{00:38:44}
\end{figure}

This subject is also called continuum mechanics.  Depending on the degrees of
freedom, the same mathematics can describe fluid flow, elastic displacement,
or another continuous medium.

\begin{classroomqa}
\audiencequestion{Would the same equations work in a world with space but no time?}
\lecturerresponse{Yes, as a purely spatial variational problem.  Physically, this is how one searches for a static solution: assume the field has no time dependence, discard all time derivatives, and solve the remaining spatial equations.  A time-independent electric or magnetic configuration is an example. \lecturetimestamp{00:40:05}}
\end{classroomqa}

\section{Guessing the Simplest Scalar Dynamics}

\subsection{Why the signs cannot all agree}

Return to the particle kinetic term \(\tfrac12\dot\phi^2\), set its irrelevant
overall coefficient to one, and try to generalize it.  A first guess might be
\begin{equation}
\mathcal{L}_{+}
=
\frac12
\left[
\phi_t^2+\phi_x^2+\phi_y^2+\phi_z^2
\right].
\label{eq:sr4-all-plus}
\end{equation}
This is a useful wrong turn.  It treats time and space with identical signs.
The resulting equation would not have the hyperbolic structure of wave
propagation.  The correct relativistic combination has one sign for time and
the opposite sign for space.

\begin{classroomqa}
\audiencequestion{Is the all-plus expression wrong because the Lagrangian should be relativistically invariant?}
\lecturerresponse{That is the deeper answer, and it is correct.  The lecture first reaches the sign more primitively: waves distinguish evolution in time from variation in space, so time cannot enter exactly like another spatial coordinate.  Lorentz invariance will shortly turn that intuition into a precise rule. \lecturetimestamp{00:43:19}}
\end{classroomqa}

Restoring a characteristic speed \(c\), the simple scalar Lagrangian density is
\begin{equation}
\boxed{
\mathcal{L}
=
\frac12
\left[
\frac{1}{c^2}\phi_t^2
-\phi_x^2-\phi_y^2-\phi_z^2
\right]
-V(\phi)
}.
\label{eq:sr4-scalar-lagrangian}
\end{equation}
For a relativistic field, \(c\) is the speed of light.  In a continuum model it
could instead be the speed of sound, the speed of an elastic wave, or another
characteristic propagation speed.

The overall factor \(1/2\) is conventional, just as it is in
\(\tfrac12m\dot\phi^2\).  Multiplying an entire action by a nonzero constant
does not change its stationary paths; the convention fixes how the parameters
in the equations are named and normalized.  Once that normalization is chosen,
it must be used consistently.

\begin{classroomqa}
\audiencequestion{Can the spatial-gradient terms be regarded as part of the potential energy?}
\lecturerresponse{Yes.  In the kinetic-minus-potential organization, the kinetic part contains time derivatives.  Terms with no time derivatives, including spatial gradients and \(V(\phi)\), belong to the potential-energy side.  The ultimate reason for their relative sign in a relativistic theory is the Lorentzian metric. \lecturetimestamp{00:45:37}}
\end{classroomqa}

\subsection{The wave equation}

Apply \eqref{eq:sr4-field-el}.  The derivative terms are
\begin{equation}
\frac{\partial\mathcal{L}}{\partial\phi_t}
=\frac{1}{c^2}\phi_t,
\qquad
\frac{\partial\mathcal{L}}{\partial\phi_i}
=-\phi_i,
\end{equation}
while
\begin{equation}
\frac{\partial\mathcal{L}}{\partial\phi}
=-\frac{dV}{d\phi}.
\end{equation}
Therefore
\begin{equation}
\begin{aligned}
&\boxed{
\frac{1}{c^2}\frac{\partial^2\phi}{\partial t^2}
-\nabla^2\phi
+\frac{dV}{d\phi}
=0
},\\[0.5em]
\nabla^2
&=
\frac{\partial^2}{\partial x^2}
+\frac{\partial^2}{\partial y^2}\\
&\quad
+\frac{\partial^2}{\partial z^2}.
\end{aligned}
\label{eq:sr4-general-scalar-eom}
\end{equation}

If \(V=0\), this reduces to
\begin{equation}
\boxed{
\frac{1}{c^2}\phi_{tt}-\nabla^2\phi=0
},
\label{eq:sr4-wave-equation}
\end{equation}
the wave equation.  The lecture recognizes the equation but postpones solving
it.  Its solutions propagate disturbances at speed \(c\).

\begin{figure}[htbp]
\centering
\includegraphics[width=0.90\textwidth]{lecture_04_figure_08.png}
\caption{The one-dimensional form \(c^{-2}\phi_{tt}-\phi_{xx}=0\), recognized as the wave equation.}
\lectureframe{00:50:30}
\end{figure}

\subsection{A field of coupled harmonic oscillators}

Now choose a quadratic potential,
\begin{equation}
V(\phi)=\frac12\mu^2\phi^2,
\qquad
\frac{dV}{d\phi}=\mu^2\phi.
\label{eq:sr4-quadratic-potential}
\end{equation}
If \(\phi\) were the position of a particle, this would be the potential of a
harmonic oscillator with spring constant \(\mu^2\).  For the field,
\begin{equation}
\boxed{
\frac{1}{c^2}\phi_{tt}-\nabla^2\phi+\mu^2\phi=0
}.
\label{eq:sr4-klein-gordon}
\end{equation}
Without the spatial derivative this is simply acceleration plus a constant
times displacement equals zero.  The gradient couples neighboring points,
turning the oscillator into a wave-bearing continuum.  In relativistic field
theory \eqref{eq:sr4-klein-gordon} is called the massive Klein--Gordon
equation; classically, \(\mu\) is the inverse length scale that modifies the
wave propagation.

\begin{figure}[htbp]
\centering
\includegraphics[width=0.90\textwidth]{lecture_04_figure_09.png}
\caption{The quadratic potential and the resulting wave equation with the restoring term \(\mu^2\phi\).}
\lectureframe{00:52:38}
\end{figure}

\begin{classroomqa}
\audiencequestion{Why does the factor \(1/2\) in \(V=\tfrac12\mu^2\phi^2\) disappear from the equation of motion?}
\lecturerresponse{Differentiating \(\phi^2\) supplies a factor of two.  The conventional \(1/2\) is inserted precisely so that \(dV/d\phi=\mu^2\phi\), just as \(V(x)=\tfrac12kx^2\) gives the restoring force \(-kx\). \lecturetimestamp{00:53:17}}
\end{classroomqa}

The equation is linear: \(\phi\) and each of its derivatives occur only to the
first power.  Terms such as \(\phi^2\) or
\(\phi\,\partial_\mu\phi\) in the equation of motion would make the dynamics
nonlinear.

\section{Lorentz Scalars and Four-Vectors}

\subsection{The action must be frame independent}

The variational machinery alone defines many possible field theories.
Relativity imposes the additional requirement that the action have the same
value and form in every inertial frame.  For the free particle this was
achieved by integrating the invariant proper time \(d\tau\) along the
worldline.  For a field, the action is integrated over spacetime, so its local
building blocks must transform appropriately.

Under a proper Lorentz transformation the spacetime measure is invariant,
\(d^4x'=d^4x\).  It is therefore sufficient to require the Lagrangian density
to be a Lorentz scalar:
\begin{equation}
\mathcal{L}'(x')=\mathcal{L}(x).
\label{eq:sr4-lagrangian-scalar}
\end{equation}

\subsection{Scalar fields}

A scalar field obeys
\begin{equation}
\boxed{\phi'(x')=\phi(x)}.
\label{eq:sr4-scalar-transform}
\end{equation}
The coordinates \(x\) and \(x'\) label the same event in two frames.  The
functional expressions used by the observers may look different, but the
number attached to that physical event is the same.

\begin{figure}[htbp]
\centering
\includegraphics[width=0.90\textwidth]{lecture_04_figure_10.png}
\caption{The defining transformation law of a scalar field: \(\phi(x)=\phi'(x')\).}
\lectureframe{00:59:33}
\end{figure}

Temperature supplied the early intuition, although a moving thermometer in a
material medium is not a perfect relativistic example.  The definition
\eqref{eq:sr4-scalar-transform}, rather than the thermometer analogy, is what
matters.

\subsection{Four-vectors}

In units \(c=1\), a boost of speed \(v\) along the \(x\)-axis acts as
\begin{equation}
\begin{aligned}
t'&=\gamma(t-vx),\\
x'&=\gamma(x-vt),\\
y'&=y,\\
z'&=z,\\
\gamma&=\frac{1}{\sqrt{1-v^2}}.
\end{aligned}
\label{eq:sr4-x-boost}
\end{equation}

\begin{figure}[htbp]
\centering
\includegraphics[width=0.90\textwidth]{lecture_04_figure_11.png}
\caption{The Lorentz transformation for a boost along the \(x\)-axis in \(c=1\) units.}
\lectureframe{01:02:26}
\end{figure}

Any four quantities that mix under boosts in the same way form a four-vector.
Examples already encountered are \(x^\mu\), \(dx^\mu\), and the four-velocity
\(u^\mu=dx^\mu/d\tau\).  Multiplying a four-vector by a scalar gives another
four-vector.  A flowing relativistic fluid supplies a field-valued example:
its four-velocity \(u^\mu(x)\) can differ from event to event.

For the same boost as in \eqref{eq:sr4-x-boost}, a generic four-vector obeys
\begin{equation}
\begin{aligned}
A'^0&=\gamma(A^0-vA^x),\\
A'^x&=\gamma(A^x-vA^0),\\
A'^y&=A^y,\\
A'^z&=A^z.
\end{aligned}
\label{eq:sr4-vector-boost}
\end{equation}

For a generic four-vector \(A^\mu=(A^0,A^x,A^y,A^z)\), the Lorentzian norm
\begin{equation}
\boxed{
A_\mu A^\mu
=(A^0)^2-(A^x)^2-(A^y)^2-(A^z)^2
}
\label{eq:sr4-vector-norm}
\end{equation}
is a scalar.  It is invariant for the same algebraic reason that
\begin{equation}
d\tau^2=dt^2-dx^2-dy^2-dz^2
\end{equation}
is invariant.  The lecture briefly runs out of unused letters for a generic
vector before accepting \(A\), helped by the suggestion \(A\) for arrow.

\begin{figure}[htbp]
\centering
\includegraphics[width=0.90\textwidth]{lecture_04_figure_12.png}
\caption{A generic four-vector field and the Lorentzian difference of squares from which its invariant norm is formed.}
\lectureframe{01:09:06}
\end{figure}

\subsection{The derivative of a scalar}

Differentiating a scalar with respect to the four spacetime coordinates
produces the four components
\begin{equation}
\partial_\mu\phi
=
\left(
\frac{\partial\phi}{\partial t},
\frac{\partial\phi}{\partial x},
\frac{\partial\phi}{\partial y},
\frac{\partial\phi}{\partial z}
\right)
\qquad(c=1).
\end{equation}
Strictly, \(\partial_\mu\phi\) is a covector; the Minkowski metric raises its
index to make \(\partial^\mu\phi\).  This distinction is harmless in the
component argument here but fixes the contraction unambiguously:
\begin{equation}
\boxed{
\partial_\mu\phi\,\partial^\mu\phi
=
\frac{1}{c^2}\phi_t^2
-\phi_x^2-\phi_y^2-\phi_z^2
}.
\label{eq:sr4-derivative-scalar}
\end{equation}

\begin{figure}[htbp]
\centering
\includegraphics[width=0.90\textwidth]{lecture_04_figure_13.png}
\caption{The invariant difference of squares built from the first derivatives of a scalar field.}
\lectureframe{01:13:55}
\end{figure}

\begin{classroomqa}
\audiencequestion{Is this derivative combination the Laplacian?}
\lecturerresponse{No.  The displayed invariant is quadratic in first derivatives.  The Laplacian contains second spatial derivatives.  The second derivatives appear only after varying the action, in the wave equation. \lecturetimestamp{01:13:22}}
\end{classroomqa}

The sign pattern in \eqref{eq:sr4-scalar-lagrangian} is now explained rather
than guessed.  The kinetic-gradient combination is the Lorentz scalar
\(\partial_\mu\phi\,\partial^\mu\phi\).  Thus the simplest quadratic theory is
\begin{equation}
\mathcal{L}
=
\frac12\partial_\mu\phi\,\partial^\mu\phi
-\frac12\mu^2\phi^2.
\label{eq:sr4-covariant-lagrangian}
\end{equation}

\begin{figure}[htbp]
\centering
\includegraphics[width=0.90\textwidth]{lecture_04_figure_14.png}
\caption{The completed quadratic scalar-field Lagrangian and its linear field equation on the board.}
\lectureframe{01:19:12}
\end{figure}

\section{How Much Freedom Is There?}

Once the rule to build scalars is known, many Lagrangians become possible.
The field \(\phi\) is a scalar, so every ordinary function \(F(\phi)\) is also
a scalar: \(\phi^2\), \(\phi^3\), \(\cosh\phi\), and so on.  Numerical
parameters are allowed as well.  A symbol such as \(\mu\) may be inserted as a
constant coefficient; the only remaining questions are its dimensions and its
physical interpretation.

\begin{classroomqa}
\audiencequestion{Can a mass be put into the Lagrangian?}
\lecturerresponse{A constant parameter certainly can.  At this stage one should ask which object's mass it is, because \(\mu\) is simply a number in the classical field theory until its physical meaning is identified.  In the quadratic relativistic scalar theory it sets the mass or inverse-length scale. \lecturetimestamp{01:14:16}}
\end{classroomqa}

Scalars may be multiplied:
\begin{equation}
F(\phi)\,\partial_\mu\phi\,\partial^\mu\phi
\end{equation}
is Lorentz invariant.  The derivative scalar itself may be squared, producing
higher powers of first derivatives.  One may also contemplate higher
derivatives, although the lecture restricts attention to fields and their first
derivatives, the standard framework that yields second-order equations of
motion.  Higher-derivative theories require additional care because they
generically introduce extra dynamical data and possible instabilities.

\begin{classroomqa}
\audiencequestion{Is the difference-of-squares construction really the metric of spacetime?}
\lecturerresponse{Yes.  The metric is \(dt^2-dx^2-dy^2-dz^2\) in the chosen signature.  The lecture has not yet emphasized the word metric, but this is exactly the spacetime geometry that will become central in general relativity. \lecturetimestamp{01:19:39}}
\end{classroomqa}

\begin{classroomqa}
\audiencequestion{Are there objects more general than four-vectors, with components such as \(xy\) or \(tx\)?}
\lecturerresponse{Yes.  Those are tensors.  A scalar has no spacetime index, a vector has one, and a rank-two tensor has components such as \(T^{xx}\), \(T^{xy}\), and \(T^{tx}\).  Most of the present scalar-field discussion avoids them; they become unavoidable in general relativity. \lecturetimestamp{01:20:49}}
\end{classroomqa}

\begin{classroomqa}
\audiencequestion{Why discuss only boosts that mix \(t\) with \(x\)?  Could a boost point in another direction and mix the spatial axes differently?}
\lecturerresponse{Certainly.  A boost along \(y\) mixes \(t\) with \(y\) and leaves \(x,z\) unchanged.  A general Lorentz transformation can be assembled from spatial rotations and a boost along a chosen axis.  It is therefore enough to establish rotational invariance and invariance under one standard-axis boost. \lecturetimestamp{01:22:06}}
\end{classroomqa}

\begin{classroomqa}
\audiencequestion{With so many possible scalar combinations, do all the resulting Lagrangians describe meaningful physics?}
\lecturerresponse{At the classical level the freedom is immense, subject to additional requirements such as sensible signs, locality, and a well-posed evolution problem.  Far fewer arbitrary choices define well-behaved fundamental quantum theories.  Modern effective field theory allows many additional interactions below a cutoff, but their range of validity must then be stated. \lecturetimestamp{01:27:57}}
\end{classroomqa}

\begin{classroomqa}
\audiencequestion{Must a field be continuous?}
\lecturerresponse{For the smooth theories under discussion, yes.  A jump produces singular derivatives and, with a quadratic gradient term, typically infinite action and energy.  More singular configurations can be treated only with additional mathematical structure; they are not the ordinary finite-energy fields being developed here. \lecturetimestamp{01:29:29}}
\end{classroomqa}

\section{A Particle Coupled to a Scalar Field}

\subsection{The free particle revisited}

The final part of the lecture joins particles and fields at the classical
level.  This is not yet the quantum statement that particles are excitations of
fields.  Instead, assume that someone has already solved for a background
scalar field \(\phi(x,t)\), and ask how a particle moves through it.

For a free relativistic particle,
\begin{equation}
A_{\rm free}=-m\int d\tau.
\label{eq:sr4-free-particle-action}
\end{equation}
In one spatial dimension and \(c=1\),
\begin{equation}
d\tau=dt\sqrt{1-\dot x^2},
\qquad
L_{\rm free}=-m\sqrt{1-\dot x^2}.
\label{eq:sr4-free-particle-l}
\end{equation}
The minus sign is chosen so that the slow-velocity expansion gives the usual
positive Newtonian kinetic term:
\begin{equation}
L_{\rm free}
\approx -m+\frac12m\dot x^2+\cdots.
\end{equation}

\begin{figure}[htbp]
\centering
\includegraphics[width=0.90\textwidth]{lecture_04_figure_15.png}
\caption{The invariant free-particle action rewritten as a coordinate-time Lagrangian.}
\lectureframe{01:34:31}
\end{figure}

\subsection{The simplest scalar coupling}

Both \(d\tau\) and \(\phi\) are scalars.  The simplest linear coupling is
therefore
\begin{equation}
\boxed{
A
=
\int d\tau\,[-m+\phi(x,t)]
}.
\label{eq:sr4-coupled-action}
\end{equation}
This choice is not unique.  One could replace \(\phi\) by any function
\(F(\phi)\).  The linear term is enough to display the physical idea.

\begin{figure}[htbp]
\centering
\includegraphics[width=0.90\textwidth]{lecture_04_figure_16.png}
\caption{The free action supplemented by the simplest Lorentz-invariant scalar coupling.}
\lectureframe{01:37:21}
\end{figure}

\begin{classroomqa}
\audiencequestion{Does \(\phi\) have to remain inside the integral?}
\lecturerresponse{Yes.  The field is evaluated on the particle's worldline, \(\phi(x(t),t)\), and generally changes from one point of the trajectory to another.  Even if the background has no explicit time dependence, \(\phi=\phi(x)\), it still changes with time when the particle moves through space.  Only a field constant along the path could be pulled out. \lecturetimestamp{01:37:47}}
\end{classroomqa}

\begin{classroomqa}
\audiencequestion{Does the \(x\) in \(\phi(x,t)\) mean the particle's position?}
\lecturerresponse{Yes.  It is \(x=x(t)\), the position of the particle at time \(t\).  The action samples the prescribed spacetime field along that trajectory. \lecturetimestamp{01:38:08}}
\end{classroomqa}

For one-dimensional motion, \eqref{eq:sr4-coupled-action} gives
\begin{equation}
L
=
-(m-\phi(x,t))\sqrt{1-\dot x^2}.
\label{eq:sr4-coupled-lagrangian}
\end{equation}
The scalar background appears in exactly the place occupied by the rest mass:
\begin{equation}
m_{\rm eff}(x,t)=m-\phi(x,t).
\label{eq:sr4-effective-mass}
\end{equation}
If \(\phi\) settles to a nearly constant nonzero value, the particle behaves as
though its mass has shifted.  This is only a classical toy model, not the full
Higgs mechanism, but it captures the advertised taste of the Higgs idea: a
scalar background can enter particle dynamics as a mass term.

\subsection{The particle equation}

The canonical momentum is
\begin{equation}
p
=
\frac{\partial L}{\partial\dot x}
=
(m-\phi)
\frac{\dot x}{\sqrt{1-\dot x^2}}.
\label{eq:sr4-coupled-momentum}
\end{equation}

\begin{figure}[htbp]
\centering
\includegraphics[width=0.90\textwidth]{lecture_04_figure_17.png}
\caption{The scalar background replaces \(m\) by \(m-\phi(x,t)\) in the relativistic canonical momentum.}
\lectureframe{01:43:38}
\end{figure}

The position derivative of the Lagrangian is
\begin{equation}
\frac{\partial L}{\partial x}
=
\frac{\partial\phi}{\partial x}
\sqrt{1-\dot x^2}.
\end{equation}
Keeping the sign fixed by the action
\eqref{eq:sr4-coupled-action}, the Euler--Lagrange equation is therefore
\begin{equation}
\boxed{
\frac{d}{dt}
\left[
(m-\phi)
\frac{\dot x}{\sqrt{1-\dot x^2}}
\right]
=
\frac{\partial\phi}{\partial x}
\sqrt{1-\dot x^2}
}.
\label{eq:sr4-coupled-eom}
\end{equation}

\begin{figure}[htbp]
\centering
\includegraphics[width=0.90\textwidth]{lecture_04_figure_18.png}
\caption{The coupled equation being assembled from the momentum and the spatial derivative of the background field; the action fixes the final sign in \eqref{eq:sr4-coupled-eom}.}
\lectureframe{01:45:24}
\end{figure}

\begin{classroomqa}
\audiencequestion{Why does \(\partial\phi/\partial x\) retain the factor \(\sqrt{1-\dot x^2}\)?}
\lecturerresponse{Because only the factor \(\phi(x,t)\) is differentiated with respect to \(x\).  The square root multiplies it in the Lagrangian and is independent of position when \(\dot x\) is held fixed in the partial derivative, so the square root remains. \lecturetimestamp{01:45:50}}
\end{classroomqa}

Expanding the total time derivative gives
\begin{align}
&(m-\phi)
\frac{d}{dt}
\left(
\frac{\dot x}{\sqrt{1-\dot x^2}}
\right)
-
\dot\phi\,
\frac{\dot x}{\sqrt{1-\dot x^2}}
\nonumber\\
&\hspace{10em}
=
\frac{\partial\phi}{\partial x}
\sqrt{1-\dot x^2}.
\label{eq:sr4-coupled-expanded}
\end{align}
If the background is static, \(\phi=\phi(x)\), its derivative along the moving
worldline is not zero.  The chain rule gives
\begin{equation}
\boxed{
\dot\phi
=
\frac{d}{dt}\phi(x(t))
=
\frac{\partial\phi}{\partial x}\dot x
}.
\label{eq:sr4-chain-rule}
\end{equation}

\begin{classroomqa}
\audiencequestion{If \(\phi\) has no explicit time dependence, why is \(\dot\phi\) not zero?}
\lecturerresponse{Because the particle changes position.  The field is sampled at \(x(t)\), so its total derivative is \(\dot\phi=\phi_x\dot x\).  A static field means \(\partial\phi/\partial t=0\), not that every moving particle sees a constant field value. \lecturetimestamp{01:48:10}}
\end{classroomqa}

\begin{figure}[htbp]
\centering
\includegraphics[width=0.90\textwidth]{lecture_04_figure_19.png}
\caption{The late expansion of the coupled equation after using the chain rule along the worldline.}
\lectureframe{01:49:11}
\end{figure}

The exact expression is not especially pretty, and the lecture does not
pretend otherwise.  Its slow-velocity limit is nevertheless familiar.  Since
\(\sqrt{1-\dot x^2}\approx1\), the leading equation is
\begin{equation}
\frac{d}{dt}\bigl[(m-\phi)\dot x\bigr]
\approx
\frac{\partial\phi}{\partial x}.
\end{equation}
If the variation of the effective mass is also treated perturbatively, this
becomes
\begin{equation}
(m-\phi)\ddot x
\approx
\frac{\partial\phi}{\partial x}.
\label{eq:sr4-nonrel-force}
\end{equation}
Indeed the low-speed Lagrangian is
\begin{equation}
L
\approx
-m+\phi+\frac12(m-\phi)\dot x^2,
\end{equation}
so at leading order the potential energy is \(U=-\phi\) and the force is
\(-dU/dx=\partial_x\phi\).  Restoring the factors of \(c\) and recovering this
Newtonian pattern is the exercise with which the lecture closes.

The route from particles to fields is now complete.  A particle history
\(\phi(t)\) supplied the \(0+1\)-dimensional prototype; a field history
\(\phi(x)\) promoted the action to a spacetime integral; the field
Euler--Lagrange equation produced wave dynamics; Lorentz scalars selected the
relativistic sign structure; and the same invariant language allowed a scalar
background to act back on a particle.
